\chapter{Ingestor agéntico (tres carriles)}
\label{sec:ingestor-agentico}

\section{Motivación}

La ingesta heurística original —un parser por formato, compilado en
el binario— cubre los formatos conocidos y se acaba: cada formato
nuevo requiere release del daemon, y cada cambio de schema aguas
arriba (Sentinel agrega un campo, FortiOS cambia un separador)
rompe el parser en producción.

El rediseño divide la ingesta en \emph{tres carriles}
independientes, cada uno con un perfil de latencia, determinismo y
superficie de cambio distinto:

\begin{enumerate}
    \item \textbf{Fast lane (VRL compilado)}: ruta caliente,
    determinística, sin LLM, sin red. Un \code{FieldConfig} ya
    validado compila a un programa VRL; el daemon lo ejecuta para
    cada batch. Es el \emph{único} carril que toca la ingesta
    productiva.
    \item \textbf{Slow lane (LLM-as-compiler, offline)}: cuando
    llega un formato nuevo o cambia uno conocido, un LLM analiza una
    muestra y propone un \code{FieldConfig} + mapping OCSF. El
    output pasa por revisión humana en la UI agéntica; al aprobarse,
    el compilador VRL emite el programa para el fast lane. El LLM
    \emph{no} toca datos en runtime.
    \item \textbf{Learning lane (drift-as-signal)}: el fast lane
    mide, a costo casi cero, métricas estructurales por batch
    (distribución de tipos, cardinalidad de campos, ratio de DLQ).
    Cuando un batch se desvía del perfil esperado más allá de un
    umbral, se dispara re-revisión en el slow lane. El fast lane no
    se para: sigue ingestando con el mapping actual mientras la
    propuesta madura.
\end{enumerate}

\section{Flujo end-to-end}

\begin{figure}[H]
\centering
\begin{tikzpicture}[
  node distance=12mm and 20mm,
  lane/.style={draw, rounded corners=2mm, minimum width=32mm, minimum height=14mm, align=center, thick},
  fast/.style={lane, fill=green!8!white, draw=green!50!black},
  slow/.style={lane, fill=orange!8!white, draw=orange!50!black},
  learn/.style={lane, fill=blue!8!white, draw=blue!50!black},
  store/.style={draw, cylinder, shape border rotate=90, aspect=0.3, minimum width=20mm, minimum height=12mm, align=center},
  arr/.style={-Latex, thick},
  arrdotted/.style={-Latex, thick, dashed},
]
  \node[fast]  (fast)    {Fast\\VRL runtime};
  \node[slow, above=of fast] (slow) {Slow\\LLM compiler\\+ human review};
  \node[learn, below=of fast] (learn) {Learning\\drift monitor};

  \node[store, left=of fast, xshift=-4mm] (raw) {raw\\events};
  \node[store, right=of fast, xshift=4mm] (canon) {canonical\\OCSF};

  \node[store, right=of slow, xshift=4mm] (lib) {mapping\\library};

  \draw[arr] (raw) -- (fast);
  \draw[arr] (fast) -- (canon);
  \draw[arr] (slow) -- (lib);
  \draw[arrdotted] (lib) -- node[above, font=\footnotesize]{compile} (fast);
  \draw[arrdotted] (fast) -- node[left, font=\footnotesize]{metrics} (learn);
  \draw[arrdotted] (learn) -- node[right, font=\footnotesize]{re-review} (slow);
\end{tikzpicture}
\caption{Tres carriles. Flechas sólidas = ruta de datos en runtime.
Flechas punteadas = feedback y compilación offline.}
\end{figure}

El contrato entre carriles es el \code{FieldConfig} + el mapping
OCSF, ambos datos. Ningún carril llama a otro por función: el slow
lane emite un artefacto que el fast lane recoge, el learning lane
emite una señal que el slow lane consume.

\section{Fast lane: VRL compilado}

El crate \filepath{platform/vrl} toma un \code{FieldConfig} y emite
un programa VRL textual. La compilación es pura:

\begin{equation*}
\text{VRL}(f) = \text{compile}(f) \qquad
\text{mapping\_hash}(f) = \text{sha256}(\text{VRL}(f))
\end{equation*}

El \code{mapping_hash} identifica al programa y se estampa en cada
evento emitido (campo \code{provenance.mapping_hash}), cerrando el
circuito entre un evento canonical y el \code{FieldConfig} que lo
produjo. El runtime es un intérprete VRL embebido; el reemplazo por
la crate \code{vrl} upstream es follow-up de performance, no de
semántica.

Consecuencias:

\begin{itemize}
    \item Cambiar un mapping no requiere release del daemon; sólo
    recompilar el VRL y reapuntar.
    \item Un batch con el mismo \code{mapping_hash} es reproducible
    byte-a-byte: ingesta determinística.
    \item El fast lane nunca llama a un LLM. Seguridad y latencia
    predecibles.
\end{itemize}

\section{Slow lane: LLM-as-compiler offline}

El LLM se usa como \emph{compilador} de ejemplos a
\code{FieldConfig}. El prompt recibe:

\begin{itemize}
    \item Muestra del formato crudo (20-100 líneas).
    \item Catálogo de tipos de entidad y relación disponibles.
    \item Schema OCSF target.
\end{itemize}

Y emite un \code{FieldConfig} propuesto + mapping OCSF. Esta
propuesta entra a la cola de review agéntica
(\code{POST /agentic_review_list}). Un humano aprueba, edita o
rechaza. Sólo tras aprobación el compilador VRL genera el programa
de fast lane.

\textbf{Propiedades}:

\begin{itemize}
    \item El LLM jamás procesa eventos de producción. Opera sobre
    muestras explícitas.
    \item El output es datos (\code{FieldConfig}), no código. El
    código lo genera el compilador determinístico.
    \item El humano es el último filtro: mismas garantías de un
    mapping curado a mano, con el boilerplate automatizado.
    \item Backend LLM intercambiable
    (\filepath{platform/local-llm}): local GGUF o API remota según
    configuración.
\end{itemize}

\section{Learning lane: drift-as-signal}

Cada batch en fast lane emite métricas baratas:

\begin{itemize}
    \item Histograma de entity types producidos.
    \item Cardinalidad de cada campo canonical.
    \item Ratio de eventos que fueron a DLQ.
    \item Distribución de timestamps (gap, skew).
\end{itemize}

El monitor de drift
(\filepath{platform/api/src/operations/agentic_drift.rs}) compara
estas métricas contra el perfil esperado del mapping. Un Kolmogorov-Smirnov
por campo, con ventana móvil. Si el drift supera el umbral:

\begin{enumerate}
    \item Se abre una propuesta de re-review apuntando al mapping
    actual.
    \item El batch \emph{sigue} ingestando con el mapping viejo —
    ingest-or-drop, no ingest-or-stall.
    \item El slow lane retoma: el LLM ve el batch anómalo + el
    mapping actual, propone delta.
    \item Humano aprueba, nuevo \code{mapping_hash}, el fast lane
    recoge la nueva versión en el siguiente batch.
\end{enumerate}

El drift no es un error, es una señal. El sistema nunca deja de
ingestar mientras aprende.

\section{Estabilidad de cada carril}

\begin{table}[H]
\centering
\small
\begin{tabular}{@{}lccc@{}}
\toprule
 & Fast & Slow & Learning \\
\midrule
Determinismo       & \checkmark         & parcial (hash estable del output) & \checkmark \\
Latencia           & runtime (ms)       & offline (minutos)                  & runtime (por batch) \\
Depende de LLM     & nunca              & sí                                 & no \\
Puede fallar caliente & sí (datos)      & n/a                                & n/a \\
Gating humano      & ninguno            & mandatorio                         & ninguno \\
\bottomrule
\end{tabular}
\caption{Perfil de cada carril. El único carril con latencia en el
camino crítico es el fast lane; sólo el slow requiere LLM y humano.}
\end{table}

\section{Relación con el resto de la arquitectura}

El ingestor agéntico vive íntegramente en la \emph{capa platform}
(ver cap.~\ref{sec:capa-platform}). El core no cambia: sigue
recibiendo \code{Event} canonical como siempre. Los cambios están
concentrados en \filepath{platform/vrl},
\filepath{platform/mapping-library},
\filepath{platform/local-llm} y
\filepath{platform/api/src/operations/agentic*.rs}.

Esto es consistente con shearing-layers: un cambio de política de
ingesta (pasar de parser-compilado a VRL-compilado-desde-LLM) se
absorbe en la capa de cambio medio sin tocar el núcleo (motor de
grafo y matcher SIMD).

\section{Qué queda fuera de este libro}

\begin{itemize}
    \item La semántica concreta del lenguaje VRL emitido y el
    mapeo canonical OCSF están en el libro de Especificación
    (ver cap.\ \emph{Formatos de log} y cap.\ \emph{MCP tools},
    dominio \code{agentic}).
    \item La evaluación cuantitativa (latencia del fast lane,
    accuracy del LLM, tasa de drift real) vive en el Paper.
\end{itemize}
